\documentclass[12pt,a4paper]{article}

\usepackage[utf8]{inputenc}
\usepackage[T1]{fontenc}
\usepackage{amsmath,amssymb}
\usepackage{booktabs}
\usepackage{geometry}
\usepackage{hyperref}
\usepackage{graphicx}
\usepackage{enumitem}
\usepackage{float}
\usepackage{caption}

\geometry{margin=2.5cm}
\hypersetup{colorlinks=true, linkcolor=blue, urlcolor=blue}

\title{FPGA-Based MPCC for the Ultra96-V2:\\Design Considerations}
\author{}
\date{}

\begin{document}

\maketitle

\noindent This document outlines considerations for implementing Model Predictive Contouring Control (MPCC) on the Ultra96-V2 FPGA, based on our experience with Pure Pursuit and an analysis of what a more advanced controller would require.

\tableofcontents
\newpage

%======================================================================
\section{Starting Point: What We Know from Pure Pursuit}
%======================================================================

The Pure Pursuit controller was our first complete FPGA implementation and served as a proof of concept for the entire workflow. It runs on the Zynq UltraScale+ ZU3EG (part \texttt{xczu3eg-sbva484-1-i}) at 100~MHz, consumes minimal resources, and completes a computation in 92--146 clock cycles (roughly 1--1.5~\textmu s).

On the real Ultra96 hardware, end-to-end compute time (including register writes, start, and readback) measures around 2.3--3.1~\textmu s. The ROS~2 overhead surrounding each call adds another 1--2~ms, which is mostly serialization and transport---not the FPGA's concern.

Resource-wise, Pure Pursuit uses about 8~BRAM tiles (4\%), 74~DSP slices (21\%), 18,700~LUTs (27\%), and 7,400~flip-flops (5\%). This leaves plenty of room on the FPGA for a more complex controller.

Key takeaways from this first implementation:

\begin{itemize}
    \item Vitis HLS 2025.1 requires \texttt{vitis-run --mode hls --tcl} (the old \texttt{vitis\_hls} binary is deprecated and co-simulation is broken through it).
    \item Scalar AXI-Lite interfaces are more reliable for co-simulation than struct-based DMA approaches.
    \item Q16.16 fixed-point arithmetic maps cleanly to DSP slices and avoids the need for floating-point IP.
    \item The development cycle (C $\rightarrow$ synthesis $\rightarrow$ co-sim $\rightarrow$ bitstream $\rightarrow$ deploy) takes roughly a day once the workflow is established.
\end{itemize}

%======================================================================
\section{Why Move Beyond Pure Pursuit}
%======================================================================

Pure Pursuit is a geometric controller. It follows a pre-computed raceline by steering toward a look-ahead point, but it has no notion of speed optimization, constraints, or prediction. The car simply follows the path at whatever speed is externally commanded.

For competitive autonomous racing, this has clear limitations. The car cannot reason about upcoming corners, cannot adjust its speed optimally, and cannot react to obstacles through its control law. It also cannot respect physical limits like the friction circle---it just does what it is told and hopes for the best.

Model Predictive Control addresses these issues by solving an optimization problem at every control step: given where the car is now, what sequence of controls over the next $N$ time steps minimizes a cost function while respecting constraints?

MPCC goes one step further by reformulating the tracking problem. Instead of tracking a reference speed profile, MPCC maximizes progress along the path while staying within physical limits. This means the controller naturally finds the fastest safe speed for any given path geometry---no pre-computed speed profile needed.

%======================================================================
\section{The Algorithm at a High Level}
%======================================================================

An MPCC controller, in simplified terms, does the following each control cycle:

\begin{enumerate}
    \item \textbf{Receive the current vehicle state} (position, heading, velocity, yaw rate, etc.)\ from localization.
    \item \textbf{Receive $N$ reference points} describing the local path to follow, including track boundary widths.
    \item \textbf{Linearize the vehicle model} at each of the $N$ prediction points (or at the current state, depending on the formulation).
    \item \textbf{Build a quadratic program} (QP) whose solution gives the optimal control sequence.
    \item \textbf{Solve the QP} using an iterative method.
    \item \textbf{Apply the first control} from the optimal sequence and repeat.
\end{enumerate}

For real-time applications, a scheme called Real-Time Iteration (RTI) is often used: instead of solving the nonlinear program to full convergence, only one or a few iterations of a Sequential Quadratic Programming (SQP) method are performed per control cycle. Since the controller runs at high frequency, the solution is continually refined and stays close to optimal.

%======================================================================
\section{Solver Selection: Why ADMM Looks Promising for FPGA}
%======================================================================

Our current MPC implementation uses Projected Gradient Descent with per-variable Gershgorin step sizes. This works well on CPU and converges quickly (around 5 iterations typically, up to 20 in difficult cases). However, it is not ideal for FPGA because each iteration involves a dense matrix-vector multiply followed by an element-wise projection---and the two operations cannot easily overlap.

The Alternating Direction Method of Multipliers (ADMM) is an alternative that may map better to hardware:

\begin{itemize}
    \item Each iteration involves a fixed matrix-vector multiply (using a pre-factored matrix that does not change within a solve), followed by element-wise soft-thresholding and projection.
    \item The matrix-vector multiply is the bottleneck and is highly parallelizable---exactly the kind of operation where FPGAs tend to excel.
    \item ADMM typically converges in 10--30 iterations for warm-started problems, which is competitive with our current solver.
\end{itemize}

The key insight is that the matrix $K = (H + \rho I)^{-1}$ can be factored once at the start of each solve and then reused for all iterations. On FPGA, this factored matrix can be stored in partitioned BRAM, allowing parallel access across rows.

There are other solvers worth considering---for example, OSQP (which also uses ADMM internally) and active-set methods. The choice depends partly on how well the solver's inner loop maps to the available DSP and memory architecture. For the ZU3EG with its 360~DSP slices, ADMM appears to be a reasonable fit, but this is not necessarily the only viable approach.

%======================================================================
\section{FPGA Performance Estimates}
%======================================================================

These estimates are rough and based on counting operations in the algorithm. Actual synthesis results would likely differ, but they give a sense of what might be feasible.

\subsection{Problem Dimensions}

For MPCC with $N=40$ prediction steps and 3 control variables per step (steering angle, longitudinal force, and path progress variable):

\begin{itemize}
    \item Decision variables: $40 \times 3 = 120$
    \item Hessian matrix: $120 \times 120 = 14{,}400$ entries
    \item Constraints: $\sim$320 (box constraints on controls + track boundary constraints)
\end{itemize}

\subsection{Parallelism in the QP Solver}

The matrix-vector multiply $K \cdot z$ ($120 \times 120$ matrix times 120-vector) dominates compute.

On a CPU (ARM Cortex-A53 at 1.2~GHz), this requires approximately 14,400 multiply-accumulate operations executed sequentially---roughly 14--43~\textmu s depending on cache behavior and instruction scheduling.

On the FPGA, if we allocate $\sim$120 DSP slices to compute all rows in parallel, the same operation completes in about 120 clock cycles at 100~MHz---approximately 1.2~\textmu s. This suggests a potential 12--36$\times$ speedup for this specific operation.

\subsection{Overall Timing}

\begin{table}[H]
\centering
\begin{tabular}{lrr}
\toprule
\textbf{Component} & \textbf{ARM A53 (est.)} & \textbf{FPGA (est.)} \\
\midrule
Model evaluation (40 steps) & $\sim$50~\textmu s & $\sim$5~\textmu s \\
QP construction & $\sim$100~\textmu s & $\sim$100~\textmu s \\
QP solve (20 ADMM iterations) & $\sim$600~\textmu s & $\sim$48~\textmu s \\
\midrule
\textbf{Total per control cycle} & \textbf{$\sim$750~\textmu s} & \textbf{$\sim$150~\textmu s} \\
\bottomrule
\end{tabular}
\caption{Estimated compute time comparison per MPCC control cycle.}
\end{table}

These numbers are estimates. The QP construction phase is harder to parallelize (it involves building the Hessian from Jacobians, which has sequential dependencies), so it does not benefit as much from the FPGA architecture.

\subsection{Resource Usage}

Including the existing Pure Pursuit (which could remain as a fallback controller):

\begin{table}[H]
\centering
\begin{tabular}{lrrrrr}
\toprule
\textbf{Resource} & \textbf{Pure Pursuit} & \textbf{MPCC (est.)} & \textbf{Combined} & \textbf{Available} & \textbf{Usage} \\
\midrule
BRAM18K & 8 & 30--50 & 38--58 & 216 & 18--27\% \\
DSP48E2 & 74 & 120--200 & 194--274 & 360 & 54--76\% \\
LUT & 18,700 & 25,000--35,000 & 43,700--53,700 & 70,560 & 62--76\% \\
FF & 7,420 & 10,000--15,000 & 17,420--22,420 & 141,120 & 12--16\% \\
\bottomrule
\end{tabular}
\caption{Estimated FPGA resource utilization for MPCC alongside Pure Pursuit.}
\end{table}

The utilization would be in the 55--76\% range for most resources. This is within typical design margins, though timing closure may require some effort at the higher end of these estimates. Designs above $\sim$75\% LUT utilization on UltraScale+ parts can sometimes be challenging to route---but this is a concern to be addressed during implementation, not a fundamental blocker.

%======================================================================
\section{Vehicle Model Considerations}
%======================================================================

Our current MPC uses a kinematic bicycle model---four states $(x, y, \psi, v)$ with velocity treated as a direct control input. This is adequate at low speeds where tire slip is negligible, but breaks down when pushing the car toward its friction limits.

For competitive racing, a dynamic model is more appropriate. A commonly used option is the dynamic single-track (bicycle) model with six states:

\begin{itemize}
    \item Position $(x, y)$
    \item Heading $\psi$
    \item Longitudinal and lateral velocities $(v_x, v_y)$
    \item Yaw rate $\dot{\psi}$
\end{itemize}

This model captures tire slip and the coupling between longitudinal and lateral forces. The tire forces can be modeled using a simplified Pacejka (Magic Formula) model or a linear tire model with saturation, depending on how much complexity is acceptable.

For FPGA implementation, the key concern with a dynamic model is the number of trigonometric and division operations per evaluation. Each model step requires $\sin$, $\cos$, $\tan$, $\text{atan2}$, and several divisions. In our Q16.16 fixed-point framework, each trig function takes roughly 50--100 cycles using Taylor series approximation. With 40 prediction steps evaluated in parallel, this is manageable---the trig functions run concurrently across all evaluation units.

The linearization (computing Jacobian matrices $\frac{\partial f}{\partial x}$ and $\frac{\partial f}{\partial u}$) is more involved with a 6-state model but follows the same structure as our current 4-state linearization: compute partial derivatives analytically and evaluate them at the operating point. The Jacobian for a 6-state, 3-control model has $6 \times 6 + 6 \times 3 = 54$ entries per prediction step, versus $4 \times 4 + 4 \times 2 = 24$ for the kinematic model.

Whether a dynamic model is necessary depends on how fast the car actually drives. Below about 3--4~m/s, the kinematic model is adequate. In our simulation environment the car reaches up to 20~m/s, which is well into the regime where tire slip dominates and a dynamic model becomes essential for accurate prediction.

%======================================================================
\section{Track Boundary Constraints}
%======================================================================

The current MPC implementation has no constraints preventing the car from leaving the track. This is fine for initial testing but needs to be addressed for real competition.

In an MPCC formulation, the most natural approach is to work in a path-relative (Frenet) coordinate frame:

\begin{itemize}
    \item The contouring error $e_c$ measures lateral deviation from the reference path.
    \item Track boundaries become simple bounds on $e_c$: the car must stay within
    \[
        -w_{\text{left}}(\theta) \leq e_c \leq w_{\text{right}}(\theta)
    \]
\end{itemize}

Here, $w_{\text{left}}$ and $w_{\text{right}}$ are the distances from the reference path to the left and right track edges, which vary along the track. These values can be extracted from the occupancy grid map that the Jetson already uses for localization.

From the FPGA's perspective, track boundary constraints are straightforward---they add two bound constraints per prediction step (left wall, right wall), bringing the total constraint count from roughly 240 to 320 for $N=40$. Since these are still simple box constraints in the Frenet frame, the projection step in the ADMM solver does not become significantly more expensive.

The track width data (one $w_{\text{left}}$ and one $w_{\text{right}}$ per reference point) can be included in the $N$ reference points sent to the FPGA each cycle. This adds $2 \times 40 = 80$ values to the per-cycle data transfer---negligible overhead.

%======================================================================
\section{Dynamic Obstacle Handling}
%======================================================================

For head-to-head racing, the car needs to avoid other vehicles on the track. There are several approaches, and the choice affects the overall system architecture more than the FPGA design specifically.

\subsection{Planner-Based Avoidance}

In this approach, obstacle avoidance is handled by a path planner running on the Jetson (which has access to LiDAR data). The planner generates a collision-free local reference path, and the FPGA's MPCC simply tracks this path.

The workflow each cycle:
\begin{enumerate}
    \item The Jetson detects obstacles from LiDAR.
    \item A local planner (e.g., gap-follow, RRT*, or lattice planner) modifies the next $N$ reference points to avoid obstacles---essentially shifting waypoints laterally within the track boundaries.
    \item The modified reference points are sent to the Ultra96.
    \item The FPGA MPCC computes optimal controls to follow the modified reference.
\end{enumerate}

Because the MPCC formulation does not use a reference speed (it maximizes progress along the path subject to physics constraints), laterally shifted waypoints do not cause speed inconsistencies. The MPCC automatically recomputes the optimal speed for the modified path geometry. In tighter corners created by obstacle avoidance, it naturally slows down; on wider paths, it speeds up.

This approach keeps the FPGA design clean---the MPCC does not need to know about obstacles at all. The downside is that the planning and control are somewhat decoupled: the planner might generate a reference that the MPCC then discovers is dynamically infeasible, or the reference might not be globally optimal when considering the car's dynamics.

\subsection{Obstacle Constraints in the MPCC}

Alternatively, obstacle avoidance can be incorporated directly into the MPCC's constraints. For each detected obstacle, a linearized collision avoidance constraint is added at each prediction step:

\[
    \vec{n}_{j,k}^{\,T} \cdot (p_k - p_{\text{obs},j}(k)) \geq d_{\text{safe}}
\]

where $\vec{n}_{j,k}$ is the direction from the obstacle to the ego vehicle, and $d_{\text{safe}}$ is a safety distance.

This adds 2 constraints per obstacle per time step (approximating the obstacle as a two-sided constraint). For 3 obstacles over $N=40$ steps, that is 240 additional constraints, bringing the total to roughly 560.

The ADMM solver handles additional constraints without fundamental difficulty---each constraint just adds one more projection operation per iteration. However, the larger constraint set increases memory requirements and may slightly increase convergence time.

This approach gives better results because the controller reasons about obstacles and dynamics simultaneously. It can find maneuvers (like late braking followed by a sharp lane change) that a planner-then-controller pipeline might miss.

\subsection{Recommendation}

For an initial implementation, the planner-based approach is likely more practical. It separates concerns, keeps the FPGA design focused on control, and does not require the FPGA to handle variable numbers of obstacle constraints. Adding obstacle constraints inside the MPCC is a natural extension if time permits.

%======================================================================
\section{System Architecture: Jetson $\leftrightarrow$ Ultra96}
%======================================================================

The overall data flow builds on the existing architecture with some modifications for MPCC.

\subsection{Per-Cycle Communication}

Each control cycle, the Jetson sends:
\begin{itemize}
    \item Vehicle state: 6--7 values (position, heading, velocities, yaw rate)
    \item $N=40$ reference points $\times$ $(x, y, \psi, w_{\text{left}}, w_{\text{right}})$ = 200 values
    \item Total: $\sim$210 values $\times$ 4 bytes = $\sim$840 bytes
\end{itemize}

The Ultra96 returns:
\begin{itemize}
    \item Optimal steering angle, longitudinal force command
    \item Optional: predicted trajectory for visualization
    \item Total: $\sim$10--20 values = $\sim$80 bytes
\end{itemize}

For communication options:
\begin{itemize}
    \item \textbf{ROS~2 (current)}: Adds 1--2~ms of overhead per round-trip, which is significant compared to the $\sim$150~\textmu s FPGA compute time. This overhead comes from serialization, DDS discovery, and middleware processing.
    \item \textbf{Raw UDP}: Estimated 0.3--0.5~ms round-trip over Ethernet. This is roughly 4--6$\times$ less overhead than ROS~2 and leaves more of the time budget available for compute.
\end{itemize}

The communication path does not need to be all-or-nothing: ROS~2 can be used locally on the Jetson for inter-node communication (LiDAR $\rightarrow$ AMCL $\rightarrow$ planner), while the Jetson-to-Ultra96 link uses raw UDP for minimal latency. The Ultra96 ARM PS receives the UDP packet, writes registers to the FPGA via AXI-Lite, waits for completion, and sends the result back.

\subsection{AXI-Lite Register Interface}

The current Pure Pursuit design uses a single AXI-Lite control bundle with $\sim$29 registers. For MPCC, the register count increases to roughly:
\begin{itemize}
    \item Control registers (mode, start, status): $\sim$5
    \item Vehicle state: $\sim$7
    \item $N \times 5$ reference values: 200
    \item Configuration parameters: $\sim$15
    \item Output registers: $\sim$5
    \item Total: $\sim$232 registers
\end{itemize}

AXI-Lite writes take approximately 1 clock cycle each at 100~MHz. Writing 232 registers takes $\sim$2.3~\textmu s---fast enough to be negligible. Reading the result (a few registers) is similarly fast.

An alternative is to use AXI DMA or AXI Stream for bulk data transfer, which would be faster for the reference points. However, AXI-Lite is simpler, well-proven in our workflow, and fast enough for the data volumes involved.

\subsection{Fallback Controller}

The Pure Pursuit controller can coexist on the same FPGA as the MPCC. If the MPCC does not produce a result within the time budget (for example, if the QP solver does not converge), the system can fall back to Pure Pursuit using the pre-loaded raceline. This provides a hardware safety net that requires no software intervention.

%======================================================================
\section{Achievable Control Frequencies}
%======================================================================

One of the clearer advantages of the FPGA is enabling higher control frequencies than what is reliably achievable on a CPU.

\begin{table}[H]
\centering
\begin{tabular}{lrcc}
\toprule
\textbf{Control Rate} & \textbf{Budget} & \textbf{FPGA} & \textbf{ARM A53} \\
\midrule
200~Hz & 5.0~ms & Feasible (97\% margin) & Feasible (with jitter risk) \\
500~Hz & 2.0~ms & Feasible (92\% margin) & Marginal \\
1~kHz & 1.0~ms & Feasible (85\% margin) & Unlikely reliable \\
2~kHz & 0.5~ms & Feasible (70\% margin) & Not feasible \\
\bottomrule
\end{tabular}
\caption{Achievable control frequencies: FPGA vs.\ ARM Cortex-A53.}
\end{table}

The practical benefit of higher control frequency is tighter tracking at the limits of adhesion. At 20~m/s, each control step at 200~Hz covers 100~mm of travel; at 1~kHz, it covers 20~mm. When operating near the friction limit, finer corrections can maintain stability where coarser ones cannot.

Whether this translates into measurably faster lap times depends on many factors---track layout, tire characteristics, localization accuracy---but the capability is there if needed, and it is something a CPU-based implementation would struggle to match reliably.

The FPGA also offers deterministic latency: the compute time is bounded by the hardware design, with no operating system jitter, cache misses, or interrupt handling delays. For a safety-critical real-time system, this predictability can be as valuable as raw speed.

%======================================================================
\section{Development Path}
%======================================================================

A reasonable order of implementation, with each phase producing a complete, testable system:

\begin{description}[style=nextline]
    \item[Phase 1: MPC on FPGA (current algorithm, proven workflow)]
    Take the existing linear MPC with kinematic model and port it to HLS. This is the lowest-risk step since the algorithm is already in C with Q16.16 fixed-point. The main work is replacing \texttt{memset}/\texttt{memcpy} with explicit loops, bounding \texttt{while} loops in angle normalization, fixing array dimensions, and adding HLS pragmas. The Projected Gradient solver should work initially---switching to ADMM can come later.

    \item[Phase 2: Upgrade to dynamic model]
    Replace the kinematic bicycle model with a dynamic single-track model. This affects the model evaluation and linearization functions. Test first on CPU (existing test framework), then synthesize for FPGA.

    \item[Phase 3: MPCC cost formulation]
    Change the cost function from position/heading tracking to contouring/lag/progress. Add the progress variable as an additional control. This changes the QP dimensions and cost weights but does not fundamentally change the solver.

    \item[Phase 4: Track boundary constraints]
    Add Frenet-frame track width constraints. This modifies QP constraint construction and requires track width data in the reference points.

    \item[Phase 5: ADMM solver]
    Replace Projected Gradient with ADMM. This is where the major FPGA performance gains come from (parallelized matrix-vector multiply). Test extensively in co-simulation.

    \item[Phase 6: Obstacle avoidance]
    Either implement the planner-based approach (CPU-side, no FPGA changes) or add linearized obstacle constraints to the MPCC.
\end{description}

Each phase is independently testable and produces a working controller. If time runs out at any phase, the system is still functional with whatever has been completed.

%======================================================================
\section{Summary of Key Design Decisions}
%======================================================================

\begin{table}[H]
\centering
\small
\begin{tabular}{lll}
\toprule
\textbf{Decision} & \textbf{Recommendation} & \textbf{Rationale} \\
\midrule
Solver algorithm & ADMM (after initial PGD) & FPGA-friendly parallel structure \\
RTI scheme & 1--3 SQP iterations/cycle & Bounded time, high-freq.\ refinement \\
Vehicle model & Dynamic single-track & Needed for competitive cornering \\
Reference data & $N$ points per cycle & Supports dynamic replanning \\
Track constraints & Frenet-frame bounds & Simple, effective, low overhead \\
Obstacle avoidance & Planner-based initially & Separation of concerns \\
Communication & Raw UDP (Jetson--Ultra96) & Minimizes latency \\
Fallback & Pure Pursuit on same FPGA & Hardware safety net \\
Fixed-point format & Q16.16 (existing) & Proven, maps to DSPs \\
Target clock & 100~MHz & Achievable on ZU3EG \\
\bottomrule
\end{tabular}
\caption{Summary of key design decisions and their rationale.}
\end{table}

\vspace{1em}
\noindent\textit{This document reflects analysis and estimates based on our Pure Pursuit implementation experience and the structure of the existing MPC codebase. Actual synthesis results, convergence behavior, and real-world performance may differ from these projections.}

%======================================================================
\section{References and Further Reading}
%======================================================================

The following are relevant papers and resources, grouped by topic. These can serve as starting points for deeper investigation and as references in the thesis.

\subsection{MPCC for Autonomous Racing}

\begin{enumerate}[label={[\arabic*]}]
    \item A.~Liniger, A.~Domahidi, and M.~Morari, ``Optimization-based autonomous racing of 1:43 scale RC cars,'' \textit{Optimal Control Applications and Methods}, vol.~36, no.~5, pp.~628--647, 2015.

    \item A.~Liniger and J.~Lygeros, ``A Noncooperative Game Approach to Autonomous Racing,'' \textit{IEEE Transactions on Control Systems Technology}, vol.~28, no.~3, pp.~884--897, 2020.

    \item J.~Kabzan et al., ``AMZ Driverless: The Full Autonomous Racing System,'' \textit{Journal of Field Robotics}, vol.~37, no.~7, pp.~1267--1294, 2020.

    \item A.~Betz et al., ``Autonomous Vehicles on the Edge: A Survey on Autonomous Vehicle Racing,'' \textit{IEEE Open Journal of Intelligent Transportation Systems}, vol.~3, pp.~458--488, 2022.
\end{enumerate}

\subsection{Real-Time Iteration (RTI) and Embedded NMPC}

\begin{enumerate}[resume, label={[\arabic*]}]
    \item M.~Diehl, H.~G.~Bock, and J.~P.~Schl\"oder, ``A Real-Time Iteration Scheme for Nonlinear Optimization in Optimal Feedback Control,'' \textit{SIAM Journal on Control and Optimization}, vol.~43, no.~5, pp.~1714--1736, 2005.

    \item R.~Verschueren, G.~Frison, D.~Kouzoupis, J.~Frey, N.~van~Duijkeren, A.~Zanelli, B.~Nober, T.~Albin, R.~Quirynen, and M.~Diehl, ``acados---a modular open-source framework for fast embedded optimal control,'' \textit{Mathematical Programming Computation}, vol.~14, pp.~147--183, 2022.
\end{enumerate}

\subsection{ADMM and QP Solvers for Embedded Systems}

\begin{enumerate}[resume, label={[\arabic*]}]
    \item B.~Stellato, G.~Banjac, P.~Goulart, A.~Bartlett, and S.~Boyd, ``OSQP: An Operator Splitting Solver for Quadratic Programs,'' \textit{Mathematical Programming Computation}, vol.~12, pp.~637--672, 2020.

    \item S.~Boyd, N.~Parikh, E.~Chu, B.~Peleato, and J.~Eckstein, ``Distributed Optimization and Statistical Learning via the Alternating Direction Method of Multipliers,'' \textit{Foundations and Trends in Machine Learning}, vol.~3, no.~1, pp.~1--122, 2011.
\end{enumerate}

\subsection{FPGA-Accelerated MPC}

\begin{enumerate}[resume, label={[\arabic*]}]
    \item J.~L.~Jerez, P.~J.~Goulart, S.~Richter, G.~A.~Constantinides, E.~C.~Kerrigan, and M.~Morari, ``Embedded Online Optimization for Model Predictive Control at Megahertz Rates,'' \textit{IEEE Transactions on Automatic Control}, vol.~59, no.~12, pp.~3238--3251, 2014.

    \item E.~N.~Hartley, J.~L.~Jerez, A.~Suardi, J.~M.~Maciejowski, E.~C.~Kerrigan, and G.~A.~Constantinides, ``Predictive Control Using an FPGA With Application to Aircraft Control,'' \textit{IEEE Transactions on Control Systems Technology}, vol.~22, no.~3, pp.~1006--1017, 2014.

    \item J.~L.~Jerez, G.~A.~Constantinides, and E.~C.~Kerrigan, ``An FPGA Implementation of a Sparse Quadratic Programming Solver for Constrained Predictive Control,'' in \textit{Proc.\ ACM/SIGDA International Symposium on FPGAs}, pp.~209--218, 2011.

    \item G.~Bel\'eman, A.~I.~Reis, and S.~S.~Mukherjee, ``FPGA-Based Model Predictive Control for Autonomous Driving,'' in \textit{Proc.\ IEEE International Conference on Field-Programmable Technology (FPT)}, 2020.
\end{enumerate}

\subsection{Vehicle Dynamics and Tire Modeling}

\begin{enumerate}[resume, label={[\arabic*]}]
    \item R.~Rajamani, \textit{Vehicle Dynamics and Control}, 2nd ed.\ Springer, 2012.

    \item H.~B.~Pacejka, \textit{Tire and Vehicle Dynamics}, 3rd ed.\ Butterworth-Heinemann, 2012.
\end{enumerate}

\subsection{F1TENTH Platform}

\begin{enumerate}[resume, label={[\arabic*]}]
    \item M.~O'Kelly, H.~Zheng, D.~Karthik, and R.~Mangharam, ``F1TENTH: An Open-source Evaluation Environment for Continuous Control and Reinforcement Learning,'' in \textit{Proc.\ NeurIPS 2019 Competition and Demonstration Track (PMLR)}, vol.~123, pp.~77--89, 2020.

    \item F1TENTH community resources: \url{https://f1tenth.org/}
\end{enumerate}

\subsection{HLS and Xilinx Tools}

\begin{enumerate}[resume, label={[\arabic*]}]
    \item Xilinx/AMD, ``Vitis High-Level Synthesis User Guide (UG1399),'' 2025.\\
    \url{https://docs.amd.com/r/en-US/ug1399-vitis-hls}

    \item Xilinx/AMD, ``Zynq UltraScale+ MPSoC Technical Reference Manual (UG1085),'' 2025.\\
    \url{https://docs.amd.com/r/en-US/ug1085-zynq-ultrascale-trm}
\end{enumerate}

\end{document}
